% Rozdział aktywny po restrukturyzacji DT.
% Roboczy podział AF 24: etykiety historyczne pozostają niezmienione.

\section{Induktywne iloczyny tensorowe funkcji testowych}
\label{sec:dt-test-tensor}
\providecommand{\atlasref}[2]{#2}

Wykorzystujemy teorię topologicznych iloczynów tensorowych
z rozdziału~\ref{sec:tensor24}, rozkład funkcji testowych
z lematu~\ref{lem:tensor-density23} oraz twierdzenie
o jądrze Schwartza~\ref{thm:schwartz-kernel23}.
Przestrzenie rozpatrujemy nad wspólnym ciałem $\K$.

\subsection{Topologia induktywna i jej własność uniwersalna dla \texorpdfstring{$\mathcal D(\Omega)$}{D(Omega)}}

\begin{definicja}[Induktywna topologia tensorowa]
\label{def:i-topology24}
Niech $E,F$ będą przestrzeniami lokalnie wypukłymi.
\emph{Induktywną topologią tensorową} $i$ na $E\otimes F$
nazywamy najsilniejszą topologię lokalnie wypukłą,
dla której kanoniczne odwzorowanie
$E\times F\to E\otimes F$ jest osobno ciągłe.
Piszemy $E\otimes_i F$, a przez
$E\widehat\otimes_i F$ oznaczamy jego uzupełnienie.
\end{definicja}

\begin{stwierdzenie}[Rola topologii induktywnej]
\label{prop:i-vs-pi24}
Dla dowolnych przestrzeni lokalnie wypukłych Hausdorffa
$E,F$ topologia induktywna $i$ istnieje, jest Hausdorffa
i nie jest słabsza od rzutowej $\pi$.
Dla każdej przestrzeni lokalnie wypukłej Hausdorffa $G$
linearyzacja daje bijekcję między osobno ciągłymi
odwzorowaniami dwuliniowymi $E\times F\to G$
a ciągłymi odwzorowaniami liniowymi $E\otimes_i F\to G$.
Jeśli $G$ jest zupełna, te ostatnie przedłużają się
jednoznacznie do $E\widehat\otimes_i F$.
W szczególności ciągłe funkcjonały liniowe na
$E\otimes_i F$ odpowiadają dokładnie osobno ciągłym
formom dwuliniowym $E\times F\to\K$. W przypadku
funkcji testowych to właśnie topologia induktywna
odpowiada twierdzeniu o jądrze Schwartza.
\end{stwierdzenie}


\begin{proof}
Rozważmy rodzinę wszystkich seminorm $h$ na
$E\otimes F$, dla których dla ustalonych $x\in E$
i $y\in F$ seminormy $h(x\otimes\cdot)$ na $F$
i $h(\cdot\otimes y)$ na $E$ są ciągłe.
Topologia generowana przez tę rodzinę zapewnia
osobną ciągłość kanonicznego odwzorowania.
Każda seminorma ciągła w dowolnej dopuszczalnej
topologii należy do tej rodziny, więc otrzymana
topologia jest najsilniejsza. Zawiera topologię
$\pi$, gdyż jej odwzorowanie kanoniczne jest
łącznie ciągłe. Własność Hausdorffa wynika z
Lematu~\ref{lem:pi-properties24}.

Niech $b:E\times F\to G$ będzie osobno ciągłe.
Dla każdej ciągłej seminormy $s$ na $G$ funkcja
$h_s=s\circ\widetilde b$ jest seminormą należącą
do opisanej rodziny: jej ograniczenia do obu
czynników przy ustalonym drugim argumencie są
ciągłe. Stąd $\widetilde b:E\otimes_i F\to G$
jest ciągłe. Odwrotnie, złożenie ciągłego liniowego
odwzorowania z osobno ciągłym odwzorowaniem
kanonicznym jest osobno ciągłe. Jednoznaczność
wynika z własności algebraicznej. Ciągłość
jednostajna odwzorowania liniowego i zupełność
$G$ dają jednoznaczne przedłużenie na uzupełnienie.
\end{proof}

Dla osobno ciągłego odwzorowania dwuliniowego $b:E\times F\to G$ oznaczamy przez $\tau$ kanoniczne odwzorowanie dwuliniowe. Strzałka przerywana oznacza jedyne odwzorowanie wyznaczone przez własność uniwersalną.
\begin{center}
\begin{tikzpicture}[>=stealth]
 \node (a) at (0,2.0) {$E\times F$};
 \node (b) at (4.4,2.0) {$E\otimes_i F$};
 \node (c) at (2.2,0) {$G$};
 \draw[->] (a) -- node[above] {$\tau$} (b);
 \draw[->] (a) -- node[below left] {$b$} (c);
 \draw[->,dashed] (b) -- node[below right] {$\widetilde b$} (c);
\end{tikzpicture}
\end{center}

\begin{stwierdzenie}[Równość topologii induktywnej i rzutowej dla przestrzeni Frécheta]
\label{prop:dt-i-pi-frechet}
Jeżeli $E,F$ są przestrzeniami Frécheta, a $G$ jest
przestrzenią lokalnie wypukłą Hausdorffa, to każde
osobno ciągłe odwzorowanie dwuliniowe $b:E\times F\to G$
jest łącznie ciągłe. W szczególności na $E\otimes F$
topologie $i$ i $\pi$ są równe. Identyczność przedłuża
się do kanonicznego izomorfizmu
$E\widehat\otimes_i F\simeq E\widehat\otimes_\pi F$.
\end{stwierdzenie}
\begin{proof}
Ustalmy ciągłą seminormę $s$ na $G$. Niech
$(V_n)$ będzie malejącą przeliczalną bazą otwartych
absolutnie wypukłych otoczeń zera w $F$. Zbiory
\[
 A_n=\bigcap_{y\in V_n}\{x\in E:s(b(x,y))\le1\}
\]
są domknięte i absolutnie wypukłe: dla każdego $y$
odwzorowanie $x\mapsto b(x,y)$ jest liniowe i ciągłe.
Ciągłość w drugiej zmiennej daje $E=\bigcup_n A_n$.
Z twierdzenia Baire'a pewien $A_n$ ma niepuste wnętrze.
Argument z Lematu~\ref{lem:separate-joint23} pokazuje,
że $A_n$ zawiera otoczenie zera $U$ (korzystamy
z wypukłości i symetrii $A_n$).
Wobec tego $s(b(x,y))\le1$ na $U\times V_n$.
Wybieramy ciągłe seminormy $p,q$ odpowiednio na $E,F$,
których otwarte kule jednostkowe zawierają się
odpowiednio w $U,V_n$. Dla dodatnich $a>p(x)$
i $c>q(y)$ mamy $s(b(x,y))\le ac$; przechodząc
do infimum, otrzymujemy
\[
 s(b(x,y))\le p(x)q(y).
\]
Obejmuje to także przypadek znikania którejś seminormy.
Ponieważ $s$ była dowolna, $b$ jest ciągłe w $(0,0)$.
Rozwinięcie
$b(x_0+x,y_0+y)-b(x_0,y_0)
=b(x_0,y)+b(x,y_0)+b(x,y)$
i ciągłość osobna zapewniają ciągłość w każdym punkcie.

Stosujemy ten wynik do kanonicznego odwzorowania
$E\times F\to E\otimes_i F$, które jest osobno ciągłe
z definicji, a którego przeciwdziedzina jest Hausdorffa
na mocy Stwierdzenia~\ref{prop:i-vs-pi24}.
Własność uniwersalna topologii rzutowej
(Twierdzenie~\ref{thm:pi-universal24}) daje ciągłość
identyczności $E\otimes_\pi F\to E\otimes_i F$.
Odwrotna identyczność jest ciągła na mocy
Stwierdzenia~\ref{prop:i-vs-pi24}.
Izomorfizm przedłuża się wraz ze swoją odwrotnością
do uzupełnień, a ich złożenia są identycznościami
przez gęstość iloczynu algebraicznego.
\end{proof}

\begin{wniosek}[Trzy topologie dla nuklearnych przestrzeni Frécheta]
\label{cor:dt-three-tensor-frechet}
Jeśli $E,F$ są przestrzeniami Frécheta i przynajmniej
jedna z nich jest nuklearna, to na $E\otimes F$
mamy $\varepsilon=\pi=i$; ich uzupełnienia są
kanonicznie izomorficzne.
\end{wniosek}
\begin{proof}
Równość $\pi=i$ wynika z poprzedniego stwierdzenia.
Równość $\varepsilon=\pi$ wynika z
Twierdzenia~\ref{thm:nuclear-pi-eps24}, w razie potrzeby
po zamianie czynników na mocy
Stwierdzenia~\ref{prop:dt-tensor-functoriality}.
\end{proof}

\begin{twierdzenie}[Produkt tensorowy funkcji testowych]
\label{thm:D-i-product24}
Niech $d_1,d_2\ge1$ będą liczbami całkowitymi, a
$\Omega\subseteq\R^{d_1}$ i $\Lambda\subseteq\R^{d_2}$
niepustymi zbiorami otwartymi. Wówczas odwzorowanie
$f\otimes g\mapsto ((x,y)\mapsto f(x)g(y))$ przedłuża się
do kanonicznego izomorfizmu liniowo-topologicznego
\[
 \mathcal D(\Omega)\widehat\otimes_i\mathcal D(\Lambda)
 \simeq \mathcal D(\Omega\times\Lambda).
\]
\end{twierdzenie}
\begin{proof}
Najpierw dowodzimy następującego faktu. Jeżeli $G$ jest
zupełną przestrzenią lokalnie wypukłą Hausdorffa, to każde
osobno ciągłe odwzorowanie dwuliniowe o wartościach w $G$,
$b:\mathcal D(\Omega)\times\mathcal D(\Lambda)\to G$,
ma jednoznaczne ciągłe liniowe przedłużenie
$B:\mathcal D(\Omega\times\Lambda)\to G$, spełniające
$B(f\otimes g)=b(f,g)$.

Dla ustalonych zwartych $K\Subset\Omega$,
$L\Subset\Lambda$ ograniczenie $b$ do
$\mathcal D_K(\Omega)\times\mathcal D_L(\Lambda)$ jest
łącznie ciągłe na mocy Stwierdzenia~\ref{prop:dt-i-pi-frechet},
gdyż oba etapy są przestrzeniami Frécheta
(Twierdzenie~\ref{thm:DK-frechet18}).
W dowodzie tego stwierdzenia dla każdej ciągłej seminormy
$s$ na $G$ otrzymaliśmy oszacowanie iloczynem dwóch
ciągłych seminorm. Ponieważ na etapach topologie
wyznaczają rosnące rodziny seminorm $p_{K,m}$ i $p_{L,r}$,
otrzymujemy
$s(b(f,g))\le C p_{K,m}(f)p_{L,r}(g)$ dla pewnych
$m,r,C$, zależnych od $s,K,L$.

Ustalamy rozkłady jedności, sześciany i odcięcia
z dowodu Lematu~\ref{lem:tensor-density23}, niezależnie
od $\Phi$. Dla $\Phi\in\mathcal D(\Omega\times\Lambda)$
ten lemat daje skończenie wiele lokalnych szeregów
$\Phi=\sum_{i,j}\sum_n f_{ijn}\otimes g_{jn}$,
gdzie dla każdej pary $(i,j)$ wszystkie czynniki
mają nośniki w ustalonych zwartych $K_i,L_j^*$.
Nad $\R$ stosujemy rzeczywistą część rozwinięcia
w sposób opisany w tym lemacie, czyli dwa składniki
rzeczywiste dla każdego indeksu.
Definiujemy
\[
 B(\Phi)=\sum_{i,j}\sum_n b(f_{ijn},g_{jn}).
\]
Dla każdej ciągłej seminormy $s$ na $G$ szereg
wewnętrzny jest bezwzględnie sumowalny względem $s$,
na mocy oszacowania na etapach i lematu.
Sumy częściowe są więc ciągiem Cauchy'ego
w topologii $G$; zupełność zapewnia istnienie granicy.

Ustalmy zwarty $M\Subset\Omega\times\Lambda$.
Lokalna skończoność rozkładów jedności daje
skończony zbiór par $(i,j)$ mogących wystąpić
dla wszystkich $\Phi\in\mathcal D_M$.
Dla każdej pary i seminormy $s$ wybieramy
$C_{ij},m_{ij},r_{ij}$ z oszacowania dwuliniowego
na etapach $K_i,L_j^*$. Wybieramy całkowite
$t_{ij}$ z $2t_{ij}>r_{ij}+d_2$.
Oszacowania Fouriera z dowodu lematu i reguła
Leibniza dla ustalonych odcięć dają
\[
 s(b(f_{ijn},g_{jn}))\le C'_{ij}
 (1+|n|)^{r_{ij}-2t_{ij}}
 p_{M,m_{ij}+2t_{ij}}(\Phi).
\]
Pochodne $\Phi$ znikają poza $M$.
Sumując po $n$, a następnie po skończonej rodzinie
par, otrzymujemy $s(B(\Phi))\le C_s p_{M,r_s}(\Phi)$
z $r_s$ równym maksimum odpowiednich rzędów
(dla pustej rodziny $B$ jest zerowe na etapie).
Stałe nie zależą od $\Phi$. W przypadku rzeczywistym
oszacowania mają co najwyżej dodatkowy czynnik $2$.

Niezależność od przedstawienia trzeba uzasadnić bez
zakładania z góry, że obie sumy definiują funkcjonały
na całej przestrzeni. Dla każdego $\lambda\in G'$ forma
skalarna $(f,g)\mapsto\lambda(b(f,g))$ jest osobno ciągła.
Aby zastosować twierdzenie o jądrze Schwartza w jego
wersji operatorowej, zauważmy, że forma ta wyznacza
ciągłe odwzorowanie $\mathcal D(\Omega)\to
\mathcal D'(\Lambda)_\beta$: dla ograniczonego
$B\subset\mathcal D(\Lambda)$ istnieje etap
$\mathcal D_L(\Lambda)$ zawierający $B$ jako zbiór
ograniczony, a łączna ciągłość formy na
$\mathcal D_K(\Omega)\times\mathcal D_L(\Lambda)$
daje dla każdego $K\Subset\Omega$ oszacowanie
$\sup_{g\in B}|\lambda(b(f,g))|\le C_B p_{K,m}(f)$
na $\mathcal D_K(\Omega)$. Ciągłość na wszystkich
etapach jest równoważna ciągłości na $\mathcal D(\Omega)$.
Twierdzenie o jądrze Schwartza
(Twierdzenie~\ref{thm:schwartz-kernel23}) daje
jednoznaczną dystrybucję $K_\lambda$, dla której
$\lambda(b(f,g))=\langle K_\lambda,f\otimes g\rangle$.
Dla każdego rozkładu tensorowego
$\Phi=\sum_n f_n\otimes g_n$ zbieżnego w przestrzeni
testowej, dla którego $\sum_n b(f_n,g_n)$ jest
zbieżny w $G$, ciągłość $\lambda$ i $K_\lambda$ daje
$\lambda(\sum_n b(f_n,g_n))=\langle K_\lambda,\Phi\rangle$.
Prawa strona nie zależy od rozkładu. Ponieważ przestrzeń sprzężona $G'$ rozdziela punkty
$G$ na mocy twierdzenia Hahna--Banacha, suma nie zależy
od wyboru rozwinięcia. Równość
$\lambda(B(\Phi))=\langle K_\lambda,\Phi\rangle$
dla wszystkich $\lambda\in G'$ dowodzi liniowości $B$.
Na tensorze elementarnym dopuszczamy jednoskładnikowe
przedstawienie, więc $B(f\otimes g)=b(f,g)$.
Jednoznaczność ciągłego przedłużenia wynika z gęstości
algebraicznych tensorów z Lematu~\ref{lem:tensor-density23}.
Z własności uniwersalnej granicy induktywnej otrzymujemy
ciągłość $B$ na całej $\mathcal D(\Omega\times\Lambda)$.

Odwzorowanie $(f,g)\mapsto f\otimes g$ o wartościach
w $\mathcal D(\Omega\times\Lambda)$ jest osobno ciągłe:
dla ustalonego $g$ o nośniku w $L$ jego ograniczenie
do $\mathcal D_K(\Omega)$ ma wartości w
$\mathcal D_{K\times L}(\Omega\times\Lambda)$, a
$p_{K\times L,m}(f\otimes g)\le p_{K,m}(f)p_{L,m}(g)$.
Stosujemy własność uniwersalną granicy induktywnej
w pierwszym czynniku. Drugi czynnik traktujemy analogicznie.
Zupełność $\mathcal D(\Omega\times\Lambda)$ wynika z
Twierdzenia~\ref{thm:D-complete18}.
Stwierdzenie~\ref{prop:i-vs-pi24} daje więc ciągłe
odwzorowanie liniowe
$J:\mathcal D(\Omega)\widehat\otimes_i\mathcal D(\Lambda)
\to\mathcal D(\Omega\times\Lambda)$.
Do zupełnego $G=\mathcal D(\Omega)\widehat\otimes_i
\mathcal D(\Lambda)$ stosujemy wykazany fakt z
kanonicznym osobno ciągłym odwzorowaniem dwuliniowym.
Otrzymujemy ciągłe $B:\mathcal D(\Omega\times\Lambda)\to G$.
Oba złożenia $JB$ i $BJ$ są identycznościami
na tensorach elementarnych, a więc na ich liniowych
powłokach. Gęstość z Lematu~\ref{lem:tensor-density23}
oraz gęstość iloczynu algebraicznego w jego uzupełnieniu
dają $JB=I$ i $BJ=I$ na całych przestrzeniach.
To jest żądany izomorfizm.
\end{proof}

\begin{przyklad}[Forma osobno ciągła, która nie jest łącznie ciągła]
\label{ex:separate-not-joint24}
Na $\mathcal D(\R)\times\mathcal D(\R)$ definiujemy
\[
 b(f,g)=\sum_{m=1}^\infty (-1)^m f^{(m)}(0)g(m).
\]
Suma jest skończona dla ustalonej funkcji $g$, ponieważ
$\supp g$ jest zwarty. Dla ustalonej $g$ funkcjonał
$f\mapsto b(f,g)$ jest więc ciągłą kombinacją skończenie
wielu pochodnych w zerze. Dla ustalonej $f$ otrzymujemy
lokalnie skończoną sumę dystrybucji punktowych
$\sum_m(-1)^mf^{(m)}(0)\delta_m$: na każdym zwartym
etapie $\mathcal D_L$ jej moduł jest ograniczony przez
$p_{L,0}(g)\sum_{m\in L\cap\N}|f^{(m)}(0)|$.
Własność uniwersalna LF zapewnia ciągłość funkcjonału. Forma $b$ jest osobno ciągła.

Przypuśćmy, że jest także łącznie ciągła. Wtedy istnieją
ciągłe seminormy $P,Q$ na $\mathcal D(\R)$ takie, że
$|b(f,g)|\le P(f)Q(g)$ dla wszystkich $f,g$.
Wybierzmy przedział zwarty $K\Subset\R$ zawierający
otoczenie zera. Ciągłość $P$ na $\mathcal D_K(\R)$
daje pewne $r\in\N_0$ i $C>0$ takie, że
$P(f)\le Cp_{K,r}(f)$ dla $f\in\mathcal D_K(\R)$.
Dla $m>r$ wybieramy $g_m\in\mathcal D(\R)$ o nośniku
w $(m-1/3,m+1/3)$ i $g_m(m)=1$. Wtedy
$|f^{(m)}(0)|\le CQ(g_m)p_{K,r}(f)$ dla wszystkich
$f\in\mathcal D_K(\R)$. Niech $\chi\in\mathcal D_K(\R)$
będzie równa $1$ w otoczeniu zera i niech
$f_N(x)=\chi(x)\cos(Nx-m\pi/2)$, gdzie $N\ge1$.
Jest to funkcja rzeczywista, więc działa nad obydwoma
ciałami. Wówczas $f_N^{(m)}(0)=N^m$, ponieważ
$\chi=1$ w otoczeniu zera. Z reguły Leibniza
dla $k\le r$ mamy
\[
 \sup_K|f_N^{(k)}|
 \le\sum_{j=0}^k\binom{k}{j}\sup_K|\chi^{(k-j)}|N^j
 \le C_{\chi,r}N^r.
\]
Po podzieleniu nierówności przez $N^r$ otrzymujemy
$N^{m-r}\le C Q(g_m) C_{\chi,r}$ dla wszystkich $N$,
co jest niemożliwe.

Forma $b$ odpowiada dystrybucji
$\sum_{m\ge1}\partial_x^m\delta_0\otimes\delta_m$
na $\R^2$; suma jest lokalnie skończona ze względu
na drugą zmienną. Jest to jawny przykład, dla którego
własność uniwersalna topologii $i$ jest silniejsza
niż własność uniwersalna topologii $\pi$.
\end{przyklad}

\begin{wniosek}[Ścisłość porównania dla dowolnych otwartych zbiorów euklidesowych]
\label{cor:dt-strict-i-pi-open}
Niech $d_1,d_2\ge1$ i niech $\Omega\subseteq\R^{d_1}$
oraz $\Lambda\subseteq\R^{d_2}$ będą niepustymi
zbiorami otwartymi. Na iloczynie algebraicznym
$\mathcal D(\Omega)\otimes\mathcal D(\Lambda)$
zachodzi $\varepsilon=\pi<i$, gdzie ostatnia
nierówność oznacza ścisłe porównanie topologii.
\end{wniosek}
\begin{proof}
Równość $\varepsilon=\pi$ wynika z nuklearności
$\mathcal D(\Omega)$ (Twierdzenia~\ref{thm:nuclear-D23}
i~\ref{thm:nuclear-pi-eps24}).
Ustalmy $x_0\in\Omega$ i wyczerpanie zwarte
$(L_m)_{m\ge1}$ zbioru $\Lambda$ z
Lematu~\ref{lem:exhaustion17}. Indukcyjnie wybieramy
różne punkty $y_m\in\Lambda\setminus L_m$.
Na każdym kroku możemy ominąć także skończony
zbiór wcześniej wybranych punktów: niepusty
otwarty zbiór $\Lambda$ nie jest zwarty, więc
nie jest sumą $L_m$ i skończonego zbioru.
Każdy zwarty podzbiór $\Lambda$ jest zawarty
w pewnym $L_M$, a zatem zawiera co najwyżej
skończenie wiele punktów $y_m$. Ponieważ każdy
punkt $\Lambda$ ma zwarte otoczenie zawarte
w $\Lambda$, rodzina tych punktów jest lokalnie
skończona i każdy $y_m$ ma otoczenie niezawierające
żadnego innego punktu tej rodziny.

Definiujemy formę dwuliniową
\[
 b(f,g)=\sum_{m\ge1}(-1)^m
 (\partial_1^m f)(x_0)g(y_m).
\]
Dla ustalonego $g$ suma jest skończona,
a każda ewaluacja pochodnej w $x_0$ jest
ciągła na $\mathcal D(\Omega)$: na etapie
zawierającym $x_0$ ogranicza ją seminorma
rzędu $m$, a na pozostałych etapach jest zerowa.
Dla ustalonego $f$ i zwartego $L\Subset\Lambda$
mamy oszacowanie
\[
 |b(f,g)|\le p_{L,0}(g)
 \sum_{m:y_m\in L}|(\partial_1^m f)(x_0)|
 \qquad(g\in\mathcal D_L(\Lambda)).
\]
Z własności uniwersalnej granicy induktywnej
(Twierdzenie~\ref{thm:LF-universal18}) wynika
osobna ciągłość $b$.

Gdyby forma była łącznie ciągła, istniałyby
ciągłe seminormy $P,Q$ z $|b(f,g)|\le P(f)Q(g)$
na mocy oszacowania z dowodu
Twierdzenia~\ref{thm:pi-universal24}.
Wybieramy zwartą kulę $K\Subset\Omega$
o środku $x_0$ i dodatnim promieniu. Istnieją
$C>0$ i $r\in\N_0$ z $P(f)\le Cp_{K,r}(f)$
dla $f\in\mathcal D_K(\Omega)$.
Dla ustalonego $m>r$ wybieramy
$g_m\in\mathcal D(\Lambda)$ z $g_m(y_m)=1$
i nośnikiem w otoczeniu $y_m$ niezawierającym
innych punktów rodziny. Istnienie takiej funkcji,
jak również odcięcia $\chi\in\mathcal D_K(\Omega)$
równego $1$ w otoczeniu $x_0$, wynika z
Lematu~\ref{lem:cutoffs17}.
Kładziemy
\[
 f_N(x)=\chi(x)\cos\bigl(N(x_1-(x_0)_1)-m\pi/2\bigr),
 \qquad N\ge1.
\]
Mamy $(\partial_1^m f_N)(x_0)=N^m$.
Reguła Leibniza, jak w poprzednim przykładzie,
daje $p_{K,r}(f_N)\le C_{\chi,r}N^r$:
pochodne czynnika trygonometrycznego względem
pozostałych współrzędnych znikają, a pochodne
względem pierwszej współrzędnej rzędu $j\le r$
mają moduł co najwyżej $N^j$.
Zatem $N^m=|b(f_N,g_m)|\le
CC_{\chi,r}Q(g_m)N^r$, co jest sprzeczne
z $m>r$.

Forma $b$ jest osobno ciągła, ale nie jest łącznie
ciągła. Z własności uniwersalnych
(Stwierdzenie~\ref{prop:i-vs-pi24} i
Twierdzenie~\ref{thm:pi-universal24})
wynika, że odpowiadający jej funkcjonał jest
ciągły dla $i$, ale nie dla $\pi$. Dowodzi to
ścisłości porównania topologii.
\end{proof}

\begin{uwaga}[Nie wolno utożsamiać trzech topologii]
\label{rem:D-not-projective24}
Na $\mathcal D(\Omega)\otimes\mathcal D(\Lambda)$
zachodzi $\varepsilon=\pi\le i$, ponieważ
$\mathcal D(\Omega)$ jest nuklearna
(Twierdzenie~\ref{thm:nuclear-D23} i
Twierdzenie~\ref{thm:nuclear-pi-eps24}).
Druga nierówność jest ostra dla wszystkich rozważanych
niepustych zbiorów otwartych, na mocy
Wniosku~\ref{cor:dt-strict-i-pi-open}.
W szczególności
$\mathcal D(\Omega)\widehat\otimes_\pi\mathcal D(\Lambda)$
nie jest \emph{kanonicznie izomorficzna topologicznie}
z $\mathcal D(\Omega\times\Lambda)$ przez odwzorowanie
$f\otimes g\mapsto f(x)g(y)$: gdyby taki izomorfizm
istniał, forma z przykładu byłaby ograniczeniem
ciągłego funkcjonału na uzupełnieniu rzutowym,
a więc byłaby łącznie ciągła. Nuklearność nie wymusza
zatem równości $\pi=i$. Jest to powód, dla którego
w twierdzeniu o jądrze Schwartza posługiwaliśmy się
formami osobno ciągłymi, a nie wyłącznie łącznie ciągłymi.
\end{uwaga}

\subsection{Powiązanie z teorią jąder i dalsza geometria tensorowa}

\begin{wniosek}[Interpretacja jądra dystrybucyjnego]
\label{cor:kernel-i24}
Niech $d_1,d_2\ge1$ będą liczbami całkowitymi,
a $\Omega\subseteq\R^{d_1}$ i $\Lambda\subseteq\R^{d_2}$
niepustymi zbiorami otwartymi. Niech
$b:\mathcal D(\Omega)\times\mathcal D(\Lambda)\to\K$
będzie formą dwuliniową. Następujące warunki są równoważne:
\begin{enumerate}[label=\textup{(\roman*)},leftmargin=*]
\item $b$ jest osobno ciągła;
\item istnieje ciągły operator liniowy
$A:\mathcal D(\Omega)\to\mathcal D'(\Lambda)_\beta$
taki, że $b(f,g)=\langle Af,g\rangle$ dla wszystkich
$f\in\mathcal D(\Omega)$ i $g\in\mathcal D(\Lambda)$;
\item istnieje dystrybucja
$K\in\mathcal D'(\Omega\times\Lambda)$ taka, że
$b(f,g)=\langle K,f\otimes g\rangle$ dla wszystkich
$f\in\mathcal D(\Omega)$ i $g\in\mathcal D(\Lambda)$.
\end{enumerate}
Operator $A$ i dystrybucja $K$ są wyznaczone jednoznacznie.
Parowanie z dystrybucją jest liniowe w funkcji testowej,
również nad $\C$; $\beta$ oznacza silną topologię
przestrzeni sprzężonej.
\end{wniosek}
\begin{proof}
\emph{(i)$\Rightarrow$(iii).}
Przez własność uniwersalną iloczynu induktywnego
(Stwierdzenie~\ref{prop:i-vs-pi24}) forma $b$
wyznacza ciągły funkcjonał na
$\mathcal D(\Omega)\otimes_i\mathcal D(\Lambda)$.
Przedłużamy go jednoznacznie na uzupełnienie (ciało
$\K$ jest zupełne) i przenosimy na
$\mathcal D(\Omega\times\Lambda)$ przez izomorfizm
z Twierdzenia~\ref{thm:D-i-product24}.
Jednoznaczność dystrybucji wynika z gęstości
liniowej powłoki tensorów elementarnych.

\emph{(iii)$\Rightarrow$(ii).}
Stosujemy odwrotną część twierdzenia o jądrze
Schwartza (Twierdzenie~\ref{thm:schwartz-kernel23}):
$A_Kf$ jest dystrybucją zadaną wzorem
$\langle A_Kf,g\rangle=\langle K,f\otimes g\rangle$,
a $A_K$ jest ciągły dla silnej topologii.
Operator jest jednoznaczny, gdyż dystrybucje
są funkcjonałami na całej przestrzeni testowej.

\emph{(ii)$\Rightarrow$(i).}
Dla ustalonego $f$ ciągłość względem $g$ wynika
z tego, że $Af$ jest dystrybucją. Dla ustalonego
$g$ ewaluacja $T\mapsto\langle T,g\rangle$ jest
ciągła w silnej topologii, ponieważ $\{g\}$
jest zbiorem ograniczonym. Składając ją z $A$,
otrzymujemy ciągłość względem $f$.

Wszystkie trzy przyporządkowania są liniowe
i wzajemnie zgodne. Otrzymujemy w szczególności
izomorfizm algebraiczny
\[
 \bigl(\mathcal D(\Omega)\widehat\otimes_i
 \mathcal D(\Lambda)\bigr)'
 \simeq\mathcal D'(\Omega\times\Lambda).
\]
Nie przypisujemy tu przestrzeni form dwuliniowych
ani przestrzeni operatorów dodatkowej topologii:
sama bijekcja liniowa nie dowodzi zgodności
dowolnie wybranych topologii na tych przestrzeniach.
\end{proof}

\begin{przyklad}[Operator o jądrze na przekątnej]
\label{ex:diag-tensor24}
Niech $d\ge1$, a $\Omega\subseteq\R^d$ będzie
niepustym zbiorem otwartym. Dystrybucja
$\delta_\Delta$ na przekątnej $\Delta=\{(x,x):x\in\Omega\}$
z Przykładu~\ref{ex:diagonal-kernel23} zadaje formę
\[
 b(f,g)=\int_\Omega f(x)g(x)\,\mathrm d\lambda_d(x),
 \qquad f,g\in\mathcal D(\Omega).
\]
Odpowiadający jej operator to regularne włożenie
$f\mapsto T_f$ do $\mathcal D'(\Omega)_\beta$,
a nie operator o wartościach w funkcjach zadany
przez całkowanie względem zwykłego jądra funkcyjnego.
W szczególności $\supp\delta_\Delta=\Delta$,
a $\delta_\Delta$ nie jest dystrybucją regularną.

W tym konkretnym przykładzie forma jest nawet
łącznie ciągła: seminorma
$p(f)=\|f\|_{L^2(\Omega)}$ jest ciągła na
$\mathcal D(\Omega)$, ponieważ na każdym zwartym
etapie zachodzi
$p(f)\le\lambda_d(K)^{1/2}p_{K,0}(f)$.
Stosujemy własność uniwersalną granicy induktywnej.
Nierówność Cauchy'ego--Schwarza daje
$|b(f,g)|\le p(f)p(g)$ także nad $\C$
(bez wprowadzania sprzężenia do definicji formy).
Zatem $b$ wyznacza ciągły funkcjonał zarówno na
uzupełnionym iloczynie rzutowym, jak i induktywnym.
Nie wynika stąd równość tych topologii: przykład
\ref{ex:separate-not-joint24} rozróżnia ich
przestrzenie ciągłych funkcjonałów.
Osobliwość jądra nie oznacza braku łącznej
ciągłości odpowiadającej mu formy.
\end{przyklad}

\begin{uwaga}[Przyszłe rozszerzenia]
\label{rem:tensor-next24}
Topologiczne produkty tensorowe pojawią się ponownie
przy operatorach zwartych i nuklearnych, własności
aproksymacji oraz algebrach operatorowych.
W szczególności normy $\pi$ i $\varepsilon$ są dwiema
różnymi normami tensorowymi już na iloczynach dwóch
nieskończenie wymiarowych przestrzeni Hilberta.
Nie rozwijamy tu zastosowań do rozwiązywania równań
różniczkowych, które należą do odrębnego skryptu.
\end{uwaga}

% Materiał z AF 25; fragment do rozdziału o iloczynie induktywnym.
\subsection{Identyfikacje przestrzeni sprzężonych do iloczynów tensorowych}

\begin{stwierdzenie}[Zestawienie kanonicznych odwzorowań]
\label{prop:summary25}
Dla każdej liczby całkowitej $d\ge1$ i ustalonego
ciała $\K\in\{\R,\C\}$ zachodzi schemat
\[
 \begin{array}{ccccc}
 \mathcal D(\R^d)&\longrightarrow&\mathcal S(\R^d)
 &\longrightarrow&\mathcal E(\R^d)\\
 \mathcal D'(\R^d)&\longleftarrow&\mathcal S'(\R^d)
 &\longleftarrow&\mathcal E'(\R^d).
 \end{array}
\]
W pierwszym wierszu strzałki oznaczają ciągłe
różnowartościowe odwzorowania liniowe o gęstych obrazach,
w drugim zaś ich transpozycje, ciągłe i różnowartościowe
zarówno w topologiach słabych-*, jak i silnych.
Żadna ze strzałek nie jest izomorfizmem topologicznym
na obraz z topologią podprzestrzeni.

Dla dodatnich całkowitych $d_1,d_2$ oraz niepustych
zbiorów otwartych $\Omega\subseteq\R^{d_1}$,
$\Lambda\subseteq\R^{d_2}$ istnieją kanoniczne
izomorfizmy liniowo-topologiczne
\[
\begin{aligned}
 \mathcal E(\Omega)\widehat\otimes_\pi\mathcal E(\Lambda)
 &\simeq\mathcal E(\Omega\times\Lambda),\\
 \mathcal S(\R^{d_1})\widehat\otimes_\pi\mathcal S(\R^{d_2})
 &\simeq\mathcal S(\R^{d_1+d_2}),\\
 \mathcal D(\Omega)\widehat\otimes_i\mathcal D(\Lambda)
 &\simeq\mathcal D(\Omega\times\Lambda).
\end{aligned}
\]
Każdy izomorfizm przedłuża odwzorowanie
$f\otimes g\mapsto ((x,y)\mapsto f(x)g(y))$.
W pierwszych dwóch wzorach można zastąpić $\pi$
przez $\varepsilon$. W trzecim wzorze nie wolno
ogólnie zastąpić $i$ przez $\pi$, choć na algebraicznym
iloczynie tensorowym przestrzeni $\mathcal D$
zachodzi $\varepsilon=\pi$.
\end{stwierdzenie}
\begin{proof}
Pierwszy wiersz wynika z Twierdzenia~
\ref{thm:zanurzenia25}, drugi z Wniosku~
\ref{cor:dual-lancuch25}; różnicę topologii
na obrazach dualnych wykazano w Przykładzie~
\ref{ex:escaping-deltas25}. Pierwsze dwa izomorfizmy
oraz dopuszczalność $\varepsilon$ wynikają z
Twierdzenia~\ref{thm:smooth-schwartz-tensor24}.
Trzeci wynika z Twierdzenia~\ref{thm:D-i-product24};
równość $\varepsilon=\pi$ wynika z nuklearności
$\mathcal D$ (Twierdzenia~\ref{thm:nuclear-D23}
i~\ref{thm:nuclear-pi-eps24}). Nie można zastąpić
$i$ przez $\pi$ już dla $\Omega=\Lambda=\R$,
na mocy Uwagi~\ref{rem:D-not-projective24}.
\end{proof}

\begin{uwaga}[Odczytywanie schematu po stronie dualnej]
\label{rem:kernel-summary25}
Jeżeli $J:X\to Y$ jest izomorfizmem liniowo-topologicznym,
to jego transpozycja $J':Y'\to X'$ jest także
izomorfizmem dla obu topologii dualnych.
Istotnie, $J'$ jest ciągłe przez równość seminorm
silnych i słabych-* z dowodu Lematu~
\ref{lem:transpozycja25}, a jego odwrotnością jest
$(J^{-1})'$, ciągłe z tego samego powodu.
Dlatego transpozycje trzech izomorfizmów tensorowych
identyfikują przestrzenie dualne produktów
z odpowiednimi przestrzeniami dystrybucji na
iloczynie, z zachowaniem topologii słabej-* i silnej.

Konkretnie, w każdym przypadku funkcjonał $K$
na przestrzeni funkcji na iloczynie odpowiada
formie dwuliniowej
$b_K(f,g)=\langle K,f\otimes g\rangle$.
Dla $\mathcal E$ i $\mathcal S$ jest to bijekcja
liniowa między dystrybucjami z odpowiedniej
przestrzeni dualnej na iloczynie a \emph{łącznie}
ciągłymi formami dwuliniowymi na parze czynników.
Wynika to z własności uniwersalnej $\pi$
(Twierdzenie~\ref{thm:pi-universal24}) i jednoznacznego
przedłużenia funkcjonału na uzupełnienie.
Dla $\mathcal D$ analogiczna bijekcja dotyczy
\emph{osobno} ciągłych form, zgodnie z
Wnioskiem~\ref{cor:kernel-i24}.
W obu przypadkach jednoznaczność jądra wynika
z gęstości liniowej powłoki tensorów elementarnych
w odpowiednim uzupełnieniu.

Nie przypisujemy tutaj przestrzeni form dwuliniowych
dodatkowej topologii. Stwierdzenie o izomorfizmach
topologicznych dotyczy dualnych przestrzeni
uzupełnionych iloczynów tensorowych, nie dowolnie
wybranych topologii na przestrzeniach form.
\end{uwaga}

\begin{uwaga}[Zakres dalszych zastosowań]
\label{rem:future25}
Identyfikacje przestrzeni sprzężonych i twierdzenie o jądrze
opisują reprezentacje ciągłych operatorów na przestrzeniach
funkcji testowych. Zastosowania do konkretnych równań
różniczkowych rozwijamy w skrypcie RRC, a ogólną teorię
operatorów w analizie funkcjonalnej. W dalszej części
niniejszego skryptu przechodzimy do analizy harmonicznej.
\end{uwaga}

\noindent\emph{Powiązania tematyczne w Bibliotece matematycznej:} \atlasref{KAT-09::cat:ex:tensor-rep}{KAT-09: własność uniwersalna iloczynu tensorowego w języku funktorów}.
